\documentclass[conference]{IEEEtran}
\IEEEoverridecommandlockouts
\usepackage{cite}
\usepackage{amsmath,amssymb,amsfonts}
\usepackage{algorithmic}
\usepackage{graphicx}
\usepackage{textcomp}
\usepackage{xcolor}
\def\BibTeX{{\rm B\kern-.05em{\sc i\kern-.025em b}\kern-.08em
    T\kern-.1667em\lower.7ex\hbox{E}\kern-.125emX}}
\begin{document}

\title{BERT-Based Named Entity Recognition for Domain-Specific Information Extraction}

\author{
\IEEEauthorblockN{Omprakash Pugazhendhi}
\IEEEauthorblockA{
  \textit{School of Computer Science and Engineering} \\
  \textit{Vellore Institute of Technology}\\
  Chennai, India \\
  omprakash.p2021@vitstudent.ac.in
}
}

\maketitle

\begin{abstract}
Named Entity Recognition (NER) is a foundational task in natural language processing that underpins information extraction, knowledge graph construction, and question answering systems. While BERT-based models achieve strong results on standard NER benchmarks, their performance degrades significantly on domain-specific corpora where entity types and linguistic patterns differ from general-domain pre-training data. In this paper, we fine-tune BERT and compare it against spaCy custom NER and BiLSTM-CRF on a domain-specific NER dataset with entity types adapted for the target domain. We evaluate entity-level F1 across PER, ORG, LOC, and MISC categories, and conduct an ablation study on the effect of training data size and domain-specific vocabulary adaptation. Our results show that BERT fine-tuned on as few as 500 domain-specific annotated sentences achieves an entity-level F1 of 0.87, outperforming BiLSTM-CRF (0.81) and spaCy custom NER (0.76) under the same training conditions.
\end{abstract}

\begin{IEEEkeywords}
named entity recognition, BERT, NLP, information extraction, fine-tuning, BiLSTM-CRF, domain adaptation
\end{IEEEkeywords}

\section{Introduction}

NER assigns semantic labels to sequences of words in text — identifying mentions of people (PER), organizations (ORG), locations (LOC), and miscellaneous entities (MISC). It is a prerequisite for downstream tasks including relation extraction, event detection, and knowledge base population \cite{nadeau2007survey}.

Pre-trained transformer models, particularly BERT \cite{devlin2018bert}, have set state-of-the-art performance on general-domain NER benchmarks such as CoNLL-2003. However, domain-specific corpora present challenges: specialized terminology, novel entity types, and entity boundary conventions that deviate from the pre-training corpus. The question of how quickly BERT adapts to a new domain with limited labeled data is practically important given annotation cost.

We compare three NER approaches under realistic data constraints:
\begin{enumerate}
\item BERT fine-tuned with a token classification head.
\item spaCy custom NER trained from scratch on domain data.
\item BiLSTM-CRF trained on domain-specific word embeddings.
\end{enumerate}

\section{Related Work}

\subsection{BERT for NER}
\cite{devlin2018bert} demonstrated that appending a token classification head to BERT and fine-tuning achieves state-of-the-art results on CoNLL-2003. \cite{li2020survey} surveys transformer-based NER approaches. Few-shot NER with BERT has been studied in \cite{huang2020few}.

\subsection{BiLSTM-CRF}
The BiLSTM-CRF architecture \cite{lample2016neural} combines bidirectional LSTMs for context encoding with a CRF layer for globally optimal label sequence decoding. It was state-of-the-art before transformer models and remains a strong baseline.

\subsection{spaCy NER}
spaCy's NER system uses a transition-based architecture with a residual CNN for feature extraction \cite{honnibal2017natural}. It is widely used in production settings due to its speed and integration with text preprocessing pipelines.

\section{Dataset}

We use the standard CoNLL-2003 English NER dataset \cite{sang2003conll} for benchmark comparison, plus a domain-specific dataset we annotated in-house comprising 3,200 sentences. The domain-specific dataset contains entity types PER, ORG, LOC, MISC annotated in technical documentation. We evaluate under three training data regimes: 100, 500, and 3,000 annotated sentences.

\section{Methodology}

\subsection{BERT Fine-Tuning}
We use \texttt{bert-base-cased} (110M parameters) from HuggingFace Transformers \cite{wolf2020transformers}. A linear classification head projects the final hidden state of each token to the number of BIO-encoded entity labels. Fine-tuning uses AdamW ($\text{lr} = 2 \times 10^{-5}$, weight decay $0.01$), batch size 32, maximum 10 epochs with early stopping on validation F1.

\subsection{BiLSTM-CRF}
Two-layer BiLSTM (hidden size 256) with 100-dimensional pre-trained GloVe embeddings and character-level CNN embeddings (50 dimensions). CRF layer for joint label sequence scoring. Trained with SGD, lr = 0.1, gradient clipping at 5.0.

\subsection{spaCy Custom NER}
spaCy 3.5 custom NER pipeline trained from blank English model. Training uses default spaCy optimizer (Adam with cyclical learning rates). Trained for 30 iterations.

\subsection{Evaluation}
Entity-level F1 (strict span match) computed using the \texttt{seqeval} library. Micro-averaged F1 reported across all entity types; per-entity-type F1 reported for analysis.

\section{Results}

\begin{table}[htbp]
\caption{Entity-Level F1 on Domain-Specific Test Set (500 Training Sentences)}
\begin{center}
\begin{tabular}{|l|c|c|c|c|c|}
\hline
\textbf{Model} & \textbf{PER} & \textbf{ORG} & \textbf{LOC} & \textbf{MISC} & \textbf{Micro F1} \\
\hline
spaCy custom & 0.84 & 0.71 & 0.79 & 0.68 & 0.76 \\
BiLSTM-CRF & 0.88 & 0.78 & 0.83 & 0.72 & 0.81 \\
BERT (fine-tuned) & \textbf{0.93} & \textbf{0.85} & \textbf{0.89} & \textbf{0.80} & \textbf{0.87} \\
\hline
\end{tabular}
\label{tab:results}
\end{center}
\end{table}

BERT consistently outperforms both baselines across all entity types. The largest gap is on ORG (0.85 vs. 0.78 BiLSTM-CRF), where domain-specific organization naming conventions benefit from BERT's contextual embeddings. MISC is the hardest category for all models.

\subsection{Data Efficiency Ablation}

\begin{table}[htbp]
\caption{Micro F1 vs. Training Data Size}
\begin{center}
\begin{tabular}{|c|c|c|c|}
\hline
\textbf{Train Size} & \textbf{spaCy} & \textbf{BiLSTM-CRF} & \textbf{BERT} \\
\hline
100 & 0.51 & 0.62 & 0.74 \\
500 & 0.76 & 0.81 & 0.87 \\
3000 & 0.85 & 0.88 & 0.91 \\
\hline
\end{tabular}
\label{tab:ablation}
\end{center}
\end{table}

BERT's advantage is largest in the low-data regime (100 sentences: F1 0.74 vs. 0.62), demonstrating that pre-trained contextual representations are especially valuable when labeled data is scarce.

\section{Conclusion}

BERT fine-tuned on 500 domain-specific sentences achieves F1 of 0.87, outperforming BiLSTM-CRF and spaCy custom NER by 6 and 11 F1 points respectively. BERT's data efficiency advantage is most pronounced in low-resource settings. Code at \texttt{github.com/omprxkash/named-entity-recognition-nlp}.

\begin{thebibliography}{00}
\bibitem{nadeau2007survey} D. Nadeau and S. Sekine, ``A Survey of Named Entity Recognition and Classification,'' Lingvisticae Investigationes, 2007.
\bibitem{devlin2018bert} J. Devlin et al., ``BERT: Pre-training of Deep Bidirectional Transformers for Language Understanding,'' NAACL, 2019.
\bibitem{li2020survey} J. Li et al., ``A Survey on Deep Learning for Named Entity Recognition,'' IEEE TKDE, 2020.
\bibitem{huang2020few} L. Huang et al., ``Few-Shot NER with Auxiliary Sentence,'' AAAI, 2020.
\bibitem{lample2016neural} G. Lample et al., ``Neural Architectures for Named Entity Recognition,'' NAACL, 2016.
\bibitem{honnibal2017natural} M. Honnibal and I. Montani, ``spaCy 2: Natural Language Understanding with Bloom Embeddings, Convolutional Neural Networks and Incremental Parsing,'' 2017.
\bibitem{sang2003conll} E. Sang and F. De Meulder, ``Introduction to the CoNLL-2003 Shared Task: Language-Independent Named Entity Recognition,'' CoNLL, 2003.
\bibitem{wolf2020transformers} T. Wolf et al., ``Transformers: State-of-the-Art Natural Language Processing,'' EMNLP, 2020.
\end{thebibliography}

\end{document}
